Frozen probing is often used to ask what information a foundation-model
representation makes accessible, but increasing probe expressivity can confound
representation quality with the probe's capacity to fit a development dataset.
We study age decoding from frozen \REVE{} EEG representations under a sealed,
predeclared external-evaluation protocol. Four heads---a final-layer
mean-pooled linear baseline, an earlier-layer linear probe, a residual head over
rich token statistics, and a two-query rich-statistics head---were trained with
the same HBN development split and \ProbeSeeds{} optimization seeds. The frozen
checkpoints were then evaluated once on \PrimarySubjects{} primary MIPDB
subjects, yielding an exact inventory of \ExternalPredictions{} subject-level
predictions. The baseline obtained mean Pearson correlation
\BaselinePearson{}. Candidate mean paired improvements were
\LayerLinearPearsonDelta{}, \RichResidualPearsonDelta{}, and
\MultiQueryPearsonDelta{}, but every hierarchical-bootstrap interval included
zero. Consequently, no tested alternative or richer head established a stable external gain
under the conjunctive decision rule. This result does not establish equivalence;
it shows that favorable average seed comparisons were insufficient evidence of
robust generalization across both optimization seeds and external subjects.
